\section{System Design and Current Implementation}

\subsection{Versioned vehicle and dock}

The initial cross-simulator specification is
\texttt{compact-semitrailer-v1}. It is a compact research-scale configuration
anchored to the NMPC paper's reported experimental dimensions, rather than a
claim of full-size highway-vehicle fidelity. Table~\ref{tab:geometry} summarizes
the frozen nominal planar geometry.

\begin{table}[ht]
  \centering
  \caption{Nominal version-1 geometry.}
  \label{tab:geometry}
  \begin{tabular}{lr}
    \toprule
    Parameter & Value \\
    \midrule
    Tractor wheelbase & $1.90\,\mathrm{m}$ \\
    Tractor body length $\times$ width & $2.90\times1.00\,\mathrm{m}$ \\
    Rear-axle-to-hitch offset & $0.00\,\mathrm{m}$ \\
    Hitch-to-trailer-axle distance & $4.00\,\mathrm{m}$ \\
    Trailer body length $\times$ width & $5.50\times1.00\,\mathrm{m}$ \\
    Wheel radius & $0.30\,\mathrm{m}$ \\
    \bottomrule
  \end{tabular}
\end{table}

The hitch offset is signed and configurable even though the nominal vehicle is
on-axle. Trailer variants use hitch-to-axle distances of $2$, $4$, and
$5\,\mathrm{m}$. An optional full-size interactive configuration now
demonstrates the same schema without replacing the compact experimental
baseline. It uses a $5.8\,\mathrm{m}$ tractor and a
$14.0\times2.55\,\mathrm{m}$ semitrailer in a $200\times100\,\mathrm{m}$ 2D
service-center yard with 20 doors along the bottom building edge and one active
$18\,\mathrm{m}$ loading zone. The target zone is highlighted by floor color
and deliberately has no physical or collidable longitudinal side walls; its
width remains part of the success containment test. The rear boundary and
building façade outside the active door are collision geometry. Faster
teleoperation limits make manual maneuvering observable at this scale; these
are demonstration choices rather than identified physical parameters.

The nominal dock interior is $6.00\times1.50\,\mathrm{m}$. Its target is the
ground-plane projection of the trailer rear-face center at
$(0.25,0.00)\,\mathrm{m}$ in the dock frame, with zero trailer yaw. Position
tolerances are $0.10\,\mathrm{m}$ longitudinally and laterally; heading
tolerance is $3^\circ$.

\subsection{Limits and timing}

The nominal steering-angle limit is $0.50\,\mathrm{rad}$ and the selected
steering-rate limit is $0.30\,\mathrm{rad/s}$. The latter is a project design
choice and will require sensitivity analysis; it is not reported by either core
paper. The soft and hard articulation thresholds are $35^\circ$ and $45^\circ$.
These values describe the compact experimental baseline. The separate
full-size interactive service-center scenario retains the $35^\circ$ warning
but terminates at $60^\circ$ articulation.
The controller period is $0.10\,\mathrm{s}$, the implemented 2D integration step is
$0.02\,\mathrm{s}$, and the MuJoCo physics step is $0.002\,\mathrm{s}$.

\subsection{Service-center presentation frame}

The service-center task retains the dock-local kinematic state while presenting
the building along the bottom edge of a conventional yard plot. If
$(x_d,y_d,\psi_d)$ is a pose in the dock frame, the renderer applies
\begin{equation}
  x_s=100-y_d,\qquad y_s=x_d,\qquad
  \psi_s=\psi_d+\frac{\pi}{2}.
\end{equation}
This transformation is a rigid $90^\circ$ rotation followed by translation,
so it preserves distances and hitch coincidence. An articulated-state
regression test separately transforms the tractor pose, trailer pose, and hitch
point and requires both transformed body hitch locations to agree with the
transformed hitch marker to numerical precision.

\subsection{Implemented software contracts}

\implemented{} components currently include immutable state and command
objects, strict vehicle and dock configuration loaders, geometry and frame
transforms, trailer variants, semantic-site definitions, state serialization,
SHA-256 configuration fingerprints, geometry-residual diagnostics, command
supervision, and an executable conformance report. Unknown, missing,
non-finite, and cross-field-inconsistent configuration values are rejected.

The 2D and MuJoCo implementations will share semantic frames for tractor rear
axle, tractor front axle, hitch, trailer axle, trailer rear face, and docking
point. MuJoCo body origins may differ, but named sites and the adapter must
recover the shared state convention.

The versioned vehicle configuration now also freezes the complete MuJoCo
identifier contract for the tractor and trailer bodies, steering links, wheel
bodies, free joint, hitch and steering joints, passive wheel-spin joints, and
four actuators. This naming layer is part of the simulator API: it permits the
implemented adapter and parity tests to resolve model elements deterministically
without depending on XML ordering.

A first MuJoCo model now implements this contract. The MJCF for the vehicle
and the nominal dock is generated from the same YAML files that define the
planar geometry, and an adapter returns the shared planar state from the named
sites at the $10\,\mathrm{Hz}$ command boundary. Wheels contact the ground
through a thin cylinder on the wheel centre plane, because full-width rigid
cylinders scrubbed at both rims and saturated the steering actuators. On a
scripted $3\,\mathrm{m}$ low-speed maneuver the model ended within
$0.009\,\mathrm{m}$ and $0.001\,\mathrm{rad}$ of a numerical integration of
the on-axle kinematic equations. This is a model-validation check on one
scenario. Direct tests now compare the project's 2D and MuJoCo backends under
shared forward and reverse command sequences; task-event and controller
transfer checks remain planned.

\subsection{Reproducibility}

Vehicle and dock YAML files contain separate schema versions and specification
identifiers. Normalized configuration fingerprints are designed to be recorded
with future experiment manifests. A new dimension set therefore changes the
result provenance without requiring controller-interface changes.

MuJoCo-only gains, damping, armature, contact width, and solver choices are
versioned in code rather than in the version-1 YAML schema. Consequently, a
reported MuJoCo experiment must also record the source revision and resolved
tuning values; the vehicle fingerprint alone is insufficient provenance.
